\section{Experimentalists: reproducing the behaviour of the methods}
\label{sec:experimentalists}

The claims of the reading list are of two kinds: mathematical statements
about statistics under stated assumptions, and operational statements about
how detectors behave on real streams. Both kinds can be checked. This section
implements the core methods of all five works in plain Python with NumPy and
SciPy, verifies the mathematics by simulation, and stages a controlled
comparison of the detector families on a synthetic benchmark with known
ground truth. Everything is seeded and deterministic; the code is listed in
Appendix~\ref{app:code} and a re-run reproduces every number and figure in
this section exactly.

\subsection{A benchmark with ground truth}
\label{sec:exp-benchmark}

Real archives are the wrong instrument for this comparison, for a reason the
reading list itself teaches: real data carry no ground truth, so any
evaluation inherits the judgement of whoever labelled them
\citep{leigh2019}. I therefore built a synthetic series designed to look like
the data my thesis system processes, thirty days of air temperature at five
minute resolution, $N = 8640$ observations, composed of a diurnal harmonic of
amplitude 5\,K around 12\,\textcelsius, a slowly wandering synoptic component
of amplitude 3\,K, and autocorrelated sensor noise, an AR(1) process with
coefficient 0.85 and innovation standard deviation 0.12\,K. Into the clean
series I injected five fault classes with disjoint supports, chosen to match
the taxonomy of \citet{campbell2013} and \citet{leigh2019}: fourteen isolated
spikes of 3 to 8\,K in either direction; a step of $+2.5$\,K lasting 36
hours, the signature of a failed radiation shield; a six hour exact flatline,
a frozen sensor repeating one value bit for bit; a six hour dithered
flatline, the same fault with residual jitter of $\pm 0.02$\,K, a sensor
mechanically stuck while its last digit still moves; and a linear calibration
drift accumulating $+2.2$\,K over four days. Three neighbouring stations
share the regional weather but carry their own small local deviations,
offsets and noise, and remain fault free. Figure~\ref{fig:series} shows the
result.

\begin{figure}[t]
  \centering
  \includegraphics[width=0.98\textwidth]{fig_series.png}
  \caption{The synthetic benchmark. Shaded bands mark the injected segment
  faults; circles mark the fourteen injected spikes. The clean signal is a
  diurnal cycle plus synoptic wandering plus AR(1) sensor noise.}
  \label{fig:series}
\end{figure}

Scoring is deliberately simple. For spikes, a fault counts as caught when a
flag lands within one sample of it. For segment faults, I report the
percentage of faulty points flagged. False alarms are flags on clean points,
reported per day for one series, because per-day per-series is the unit an
operator actually experiences.

\subsection{Grubbs' mathematics, verified, and masking, reproduced}
\label{sec:exp-grubbs}

The first experiment checks Equation~\eqref{eq:grubbscrit} against brute
force. For each sample size $n$ from 4 to 30, I drew 20\,000 standard normal
samples, computed the extreme studentised deviate of each, and took the
empirical 95th percentile. Figure~\ref{fig:grubbs}, left panel, overlays the
Monte Carlo estimates on the closed-form curve: the largest absolute
discrepancy across all 27 sample sizes is 0.018, comfortably within Monte
Carlo error. The formula Grubbs tabulated by hand in the 1960s survives
20\,000-fold replication without visible strain, which is worth pausing on:
the paper's tables were the era's entire computational infrastructure, and
they were right.

\begin{figure}[t]
  \centering
  \includegraphics[width=0.98\textwidth]{fig_grubbs.png}
  \caption{Left: two-sided 5\,\% critical values of the Grubbs statistic,
  exact formula against Monte Carlo with 20\,000 draws per sample size.
  Right: the masking experiment. A planted pair of equal outliers defeats
  the iterated Grubbs test, which flags nothing, while the generalised ESD
  procedure recovers both.}
  \label{fig:grubbs}
\end{figure}

The second experiment reproduces the failure mode \citet{pearson1936}
predicted on paper. I planted two equal outliers at 10.6 into a tight sample
of eighteen inliers around 10.0, giving a contaminated sample of twenty. The
pair inflates the sample standard deviation enough that the largest
studentised deviate reaches only 2.56, below the two-sided critical value
of roughly 2.71 at $n = 20$, so Grubbs' test, applied one observation at a
time, flags nothing at all and stops immediately. Rosner's generalised ESD
procedure with $r = 4$, which computes all four successive statistics before
deciding, finds the first statistic insignificant and the second one large,
and correctly declares exactly two outliers \citep{rosner1983}. The masking
literature is sometimes presented as a curiosity; on this evidence it is a
two-line trap that any pipeline applying a fixed-critical sequential test
will eventually fall into.

\subsection{The robust score does what it promises}
\label{sec:exp-ih}

The same twenty-point sample separates the classical and modified z-scores
cleanly. Figure~\ref{fig:modz} plots both scores for every observation. The
planted pair never reaches the classical rule of thumb, $\lvert z \rvert$
peaks at 2.56 because the outliers corrupt both the mean and the standard
deviation used to judge them, while the modified z-score of
Equation~\eqref{eq:modz}, whose median and MAD ignore the pair almost
entirely, assigns them $\lvert M \rvert = 5.17$, well beyond the 3.5 rule of
\citet{iglewicz1993}. On its home ground, a single contaminated batch, the
booklet's recommendation behaves exactly as advertised, and the experiment
needed eleven lines of code. The interesting question is what happens when
the score leaves home, which is the next experiment.

\begin{figure}[t]
  \centering
  \includegraphics[width=0.98\textwidth]{fig_modz.png}
  \caption{Classical and modified z-scores on the masking sample. The robust
  score isolates the planted pair; the classical score never clears its own
  threshold.}
  \label{fig:modz}
\end{figure}

\subsection{The threshold does not transfer}
\label{sec:exp-threshold}

Operational spike detectors, including the one in the system I worked on,
apply the modified z-score not to a batch but to a short window sliding
along the stream, and they run it with thresholds far above 3.5, values
like 30 at five minute resolution or 300 at hourly resolution appear in
production configurations. The literature gives no derivation for such
values, so I measured what the choice actually trades. For centred windows of
7, 25 and 97 samples, I swept the threshold from 2 to 400 and recorded, on
the benchmark, the false alarm rate on clean points and the recall of the
fourteen injected spikes. Figure~\ref{fig:thresholds} shows both curves.

\begin{figure}[t]
  \centering
  \includegraphics[width=0.95\textwidth]{fig_thresholds.png}
  \caption{Rolling modified z-score on the benchmark: false alarms per day
  (top) and spike recall (bottom) as functions of the threshold, for three
  window lengths. The dashed line marks the textbook 3.5; the shaded band
  marks the range of thresholds found in operational configurations.}
  \label{fig:thresholds}
\end{figure}

Three results stand out. First, the textbook constant is noisy here: at
window 7 the 3.5 rule catches all fourteen spikes but pays 1.63 false alarms
per day on a single clean series. That sounds modest until it meets network
arithmetic: across a hundred series it is over 160 nuisance flags a day
landing on whoever reads the reports, precisely Anscombe's premium, now
payable in reviewer attention. Second, the operational values silence the
false alarms and quietly give up detection: at threshold 30 the same
detector catches 9 of 14 spikes, 64\,\%, and the misses are not tail cases
but ordinary 3 to 5\,K spikes. Third, and most useful, no setting wins. At
window 7 full recall survives only up to a threshold of about 4, while
false alarms persist until about 18; at window 25 the two boundaries are
2.5 and 5; at window 97 the detector cannot reach full recall at any
threshold, topping out at 64\,\%, because an eight hour window swallows so
much diurnal variation that its MAD dwarfs a moderate spike. There is no
clean operating window anywhere in the plane I searched: on this benchmark
the statistic itself, a robust score built for exchangeable batches and
applied to an autocorrelated stream, cannot deliver both silence and full
detection, and no threshold can be tuned into fixing that. The practical
moral is sharper than the usual advice to choose thresholds carefully: the
window, the threshold, the sampling interval and the noise scale form one
coupled design problem, and a constant carried from a 1993 booklet, or from
another network's configuration file, is a measurement about someone else's
data.

\subsection{The battery against the fault classes}
\label{sec:exp-battery}

The final experiment runs the whole lineage side by side: the two global
batch methods, the rolling modified z-score at both the textbook and an
operational threshold, and one representative of each family in
Equations~\eqref{eq:range} to~\eqref{eq:regression}, a step test at 3\,K per
five minutes, persistence tests in both the naive equality form and the
variance form of Equation~\eqref{eq:persistence}, a spatial test flagging
departures beyond three residual standard deviations from the offset-adjusted
neighbour median, and a one-step-ahead regression detector in the spirit of
Equation~\eqref{eq:regression}, forecasting each value from its predecessor
plus a climatological time-of-day increment learned from the first five
days. Figure~\ref{fig:detectors} reports the full matrix of recall by fault
class alongside each detector's false alarm rate.

\begin{figure}[t]
  \centering
  \includegraphics[width=0.98\textwidth]{fig_detectors.png}
  \caption{Share of faulty points flagged, by detector and fault class, with
  false alarms per day in the right margin. Every row is one implemented
  detector run over the identical benchmark.}
  \label{fig:detectors}
\end{figure}

The matrix compresses several of the reading list's central claims into one
picture. The two global batch methods, Grubbs' test applied to the month as
if it were a sample and the modified z-score used the same way, flag nothing
whatsoever: not one of the fourteen spikes, some of them 8\,K excursions,
because a month of diurnal cycling inflates the scale estimate until nothing
looks surprising. This is the cleanest possible demonstration that the
assumption, exchangeability, dominates the formula. The step test catches
every spike cheaply but sees a level shift only at its two edges, flagging
2 of 432 points of the step segment, and is blind to everything slow or
still, exactly the profile its 1990s designers understood. The persistence
pair stages the report's neatest lineage result: run detection by exact
equality catches the bit-for-bit flatline perfectly and the dithered
flatline not at all, while the variance form used by the Mesonet,
Equation~\eqref{eq:persistence}, catches both without a single false alarm.
A sensor that is stuck but still twitching in its last digit defeats the
naive implementation and not the 2000 design; the fix was in the reading
list all along, a point with direct relevance to any modern reimplementation
that reaches for \texttt{value != previous} out of habit.

The spatial test is the matrix's quiet winner: every spike, essentially the
whole step, four fifths of both flatlines, and, uniquely in the battery,
one third of the drift, all at 0.1 false alarms per day. It sees the drift
only from the point the accumulated offset crosses its band, roughly two
days in, which is the honest shape of spatial detection: redundancy buys
sensitivity to exactly the faults that leave a single series looking
internally plausible, at the price of latency proportional to the fault's
slowness. The one-step regression detector catches all spikes and the step's
entry, and nothing of the drift, reproducing in miniature the type-resolved
finding of \citet{leigh2019} that one-step-ahead surprise is the wrong
instrument for slow degradation, since a drifting sensor surprises no one
from one observation to the next. And the rolling modified z-score rows
restate Section~\ref{sec:exp-threshold} in company: at 3.5 it is the most
trigger-happy row of the battery, while at 30 it misses over a third of the
spikes, faults the far simpler step test catches in full.

No single row dominates, so the last row of Figure~\ref{fig:detectors}
measures the obvious composition: the union of the step test, the variance
persistence test and the spatial test. It catches every spike, the entire
step and both flatlines completely, retains the spatial test's one third of
the drift, and costs 0.6 false alarms per day, less than any competitive
single detector achieves on its own. A battery of three cheap, transparent
rules from the 1990s outperforms every individual method in the table,
including the statistically fancier ones. That is, in one figure, why
\citet{shafer2000} shipped a battery rather than a test, why
\citet{campbell2013} organise checks as a sequential pipeline, and why
\citet{leigh2019} end by recommending ensembles; the only fault class the
union leaves substantially uncovered, slow drift, is exactly the class the
whole literature agrees is hardest, and the case for the model-based methods
of the current era rests on it.

\subsection{What this exercise cannot show}
\label{sec:exp-limits}

The benchmark is kind to every detector in it. Its clean process is
stationary in a way no real month of weather is, its faults are crisp,
disjoint and generously sized, there are no missing values, no fouling that
mimics real signals, no storms that mimic faults, and its neighbour stations
agree more politely than real terrain ever allows. The absolute numbers
above, 1.63 false alarms per day, 64\,\% recall, 0.018 agreement, are
therefore properties of the benchmark, not of nature, and only the ordering
of methods and the shape of the trade-offs should travel. I also set the
battery's constants myself, the 3\,K step limit, the 0.02\,K variance floor,
the three sigma spatial band, by exactly the kind of judgement this report
criticises elsewhere, which is worth admitting: within an afternoon of
experiment design I had reproduced not only the methods of the reading list
but its epistemic situation, thresholds chosen by a person who believed they
were reasonable. Finally, the feature-based branch of \citet{talagala2019}
is discussed but not implemented here; its extreme value machinery deserves
a fuller treatment than a synthetic benchmark with fourteen events can
support.
